\section{Giới thiệu và xác định bài toán}
\label{sec:gioi-thieu}

\subsection{Bối cảnh}
\label{sec:boi-canh}

Phân biệt chính xác giữa mưa (Rain) và tuyết (Snow) có ý nghĩa thực tiễn trực tiếp đối
với cảnh báo giao thông đường bộ, lập kế hoạch nông nghiệp và vận hành lưới điện/hạ tầng
đô thị ở các vùng khí hậu lạnh, vì hai pha giáng thủy này gây ra hệ quả vận hành khác
nhau dù xảy ra trong cùng một dải nhiệt độ. Trong thực tế, các trạm quan trắc tự động
không cung cấp chuỗi dữ liệu liên tục và đầy đủ: cảm biến hỏng, mất điện, lỗi truyền tín
hiệu hoặc bảo trì định kỳ khiến một phần quan trắc bị thiếu tại các thời điểm ngẫu nhiên
hoặc theo từng đoạn liên tục. Một mô hình phân loại chỉ được kiểm chứng trên dữ liệu đầy
đủ có thể suy giảm hiệu năng mạnh khi gặp dữ liệu thiếu trong vận hành thực tế, nhưng mức
độ suy giảm theo từng họ phương pháp (học máy cổ điển, biến đổi chuỗi thời gian, học sâu)
hiện chưa được so sánh một cách hệ thống trên cùng một bộ dữ liệu và cùng một giao thức
giả lập thiếu dữ liệu.

\subsection{Phát biểu bài toán}
\label{sec:phat-bieu-bai-toan}

Bài toán được phát biểu như một bài toán phân loại nhị phân trên chuỗi thời gian đa biến:

\begin{itemize}
  \item \textbf{Đầu vào:} cửa sổ trượt $X \in \mathbb{R}^{24 \times 5}$, tương ứng với
        24 giờ quan trắc liên tiếp của 5 biến khí tượng (nhiệt độ, độ ẩm, tốc độ gió, áp
        suất, và tầm nhìn).
  \item \textbf{Đầu ra:} nhãn nhị phân $y \in \{\text{Rain}, \text{Snow}\}$ tại thời
        điểm $t+1$, ngay sau cửa sổ quan trắc.
  \item \textbf{Điều kiện ràng buộc:} tại thời điểm suy luận, một phần của $X$ có thể bị
        thiếu (giả lập theo mẫu điểm, đoạn liên tiếp hoặc khối), và mô hình phải đưa ra
        dự đoán mà không được truy cập nhãn hoặc dữ liệu tương lai.
\end{itemize}

\subsection{Câu hỏi nghiên cứu}
\label{sec:cau-hoi-nghien-cuu}

Câu hỏi nghiên cứu trung tâm của đồ án là: \emph{trong số các họ phương pháp học máy cổ
điển dựa trên đặc trưng, họ biến đổi ROCKET, và học sâu chuyên biệt cho dữ liệu thiếu,
phương pháp nào giữ được hiệu năng phân loại Rain/Snow tốt nhất khi tỷ lệ và kiểu mất dữ
liệu tăng lên?} Các câu hỏi phụ bao gồm: mức suy giảm hiệu năng có khác nhau theo kiểu
mất dữ liệu (điểm/đoạn/khối) hay không, và việc đưa thêm thông tin về vị trí mất dữ liệu
(mask, $\Delta t$) có cải thiện khả năng chống chịu của từng họ phương pháp hay không.

\subsection{Đóng góp}
\label{sec:dong-gop}

\begin{itemize}
  \item Xây dựng một giao thức benchmark thống nhất trên bộ dữ liệu Szeged, trong đó ba
        kiểu mất dữ liệu (điểm, đoạn liên tiếp, khối) ở các tỷ lệ 10/30/50\% được giả lập
        bằng một bộ mặt nạ dùng chung cho mọi phương pháp, đảm bảo so sánh công bằng.
  \item Triển khai pipeline cho ba hướng phương pháp -- Classic ML, ROCKET
        và học sâu -- trên cùng giao thức. Phần thực nghiệm hiện có hoàn tất
        năm mô hình DL: GRU-Fill, GRU+mask/$\Delta t$, GRU-D, BRITS và SAITS;
        chưa kết luận so sánh với hai nhánh còn lại hoặc UniTS.
  \item Phân tích định lượng Macro F1, PR-AUC, chỉ số từng lớp và ablation
        mask/$\Delta t$ trên 150 lượt DL; trình bày rõ các phân tích thống kê,
        tuning, mùa/nhiệt độ và giải thích mô hình chưa thực hiện.
  \item Lưu cấu hình, môi trường, checksum, checkpoint và prediction theo
        từng case/seed; cung cấp bảng/hình tái tạo được để đối chiếu kết quả
        và đánh đổi hiệu năng--chi phí.
\end{itemize}

\subsection{Cấu trúc báo cáo}
\label{sec:cau-truc-bao-cao}

Phần còn lại của báo cáo được tổ chức như sau. Chương~\ref{sec:du-lieu} trình bày bộ dữ
liệu Szeged và các vấn đề chất lượng dữ liệu liên quan. Chương~\ref{sec:tong-quan-tai-lieu}
tổng quan các hướng tiếp cận liên quan trong tài liệu. Chương~\ref{sec:phuong-phap-de-xuat}
trình bày pipeline và các phương pháp được triển khai. Chương~\ref{sec:thiet-lap-thuc-nghiem}
mô tả thiết lập thực nghiệm. Chương~\ref{sec:ket-qua-va-phan-tich} trình bày và phân tích
kết quả. Chương~\ref{sec:ket-luan} tổng kết và đề xuất hướng phát triển.
